\documentclass[11pt,a4paper]{article}
\usepackage{nabiadoc}

\title{Gaussian Process Forecasting\\[0.3em]\large Kernels, posterior prediction, marginal likelihood and the \texttt{nabia} implementation}
\author{nabia documentation}
\date{October 2026}

\begin{document}
\maketitle

\begin{abstract}
This note documents \code{GaussianProcess} and the kernels in \code{Kernels.h}.
We define Gaussian process regression on the time index and prove that the
\code{nabia} kernels are valid covariance functions. The proofs cover the squared exponential kernel via
its power series, the periodic kernel via a warping of time onto a circle, and sums and
products via the Schur product theorem. We derive the gradient of every kernel with
respect to its log-hyperparameters and the posterior predictive distribution.
We also derive the log marginal likelihood and its gradient, using Jacobi's formula for
the log-determinant. Further results: how normalizing the data changes the likelihood, the
long-horizon behavior of forecasts (mean reversion for decaying kernels,
exact periodicity for the periodic kernel), and the Cholesky-based algorithms
with their costs. The last section maps everything to the code.
\end{abstract}

\tableofcontents

%==============================================================================
\section{Model}

\begin{definition}
A \emph{Gaussian process} $f\sim\mathcal{GP}(0,k)$ on $\R$ is a random
function such that, for any finite set of times $t_1,\dots,t_m$, the vector
$(f(t_1),\dots,f(t_m))$ is Gaussian with mean zero and covariance
$K_{ij}=k(t_i,t_j)$. A function $k$ can serve as a covariance function iff it is
symmetric and \emph{positive semi-definite}: every such Gram matrix $K$ is positive
semi-definite.
\end{definition}

\code{nabia} places the observations of a series at $t=0,1,\dots,T-1$ and
models
\begin{equation}\label{eq:model}
  y_t = m + s\,\big(f(t) + \varepsilon_t\big),\qquad f\sim\mathcal{GP}(0,k),\qquad
  \varepsilon_t\overset{\text{iid}}{\sim}\N(0,\sigma_n^2),
\end{equation}
where $m$ is the sample mean and $s$ the sample standard deviation (or $s=1$
without normalization, Section~\ref{sec:normalization}). We work with the
standardized series $z_t=(y_t-m)/s$, so that $z = f + \varepsilon$ with
$z\sim\N(0,K)$, $K = K_f + \sigma_n^2 I$, and $(K_f)_{ij}=k(t_i,t_j)$.

The kernel encodes what the series looks like. Its variance sets the amplitude,
its length scale sets how fast the function can change, and a period sets which
patterns repeat. The data then choose these hyperparameters by maximizing the
marginal likelihood (Section~\ref{sec:ml}).

%==============================================================================
\section{Kernels}

All \code{nabia} kernels are stationary: they depend on $\tau=t-t'$ only.
Their hyperparameters $\theta_j>0$ are exposed as $\log\theta_j$ (see
Section~\ref{sec:optimization}), and \code{gradient()} returns
$\partial k/\partial\log\theta_j=\theta_j\,\partial k/\partial\theta_j$.

\subsection{Definitions and log-parameter gradients}

\paragraph{Squared exponential} (\code{SquaredExponentialKernel}, hyperparameters $\sigma^2,\ell$):
\[
  k(\tau)=\sigma^2\exp\!\Big(-\frac{\tau^2}{2\ell^2}\Big),\qquad
  \frac{\partial k}{\partial\log\sigma^2}=k,\qquad
  \frac{\partial k}{\partial\log\ell}=k\,\frac{\tau^2}{\ell^2}.
\]
The second derivative follows from $\partial(\tau^2/2\ell^2)/\partial\ell=-\tau^2/\ell^3$,
multiplied by $\ell$. Sample paths are infinitely differentiable.

\paragraph{Matérn 5/2} (\code{Matern52Kernel}, $\sigma^2,\ell$): with $s=\sqrt5|\tau|/\ell$,
\[
  k=\sigma^2\Big(1+s+\frac{s^2}{3}\Big)e^{-s},\qquad
  \frac{\partial k}{\partial\log\sigma^2}=k,\qquad
  \frac{\partial k}{\partial\log\ell}=\sigma^2\,\frac{s^2(1+s)}{3}\,e^{-s}.
\]
\emph{Derivation.} $\frac{dk}{ds}=\sigma^2e^{-s}\big[(1+\tfrac{2s}{3})-(1+s+\tfrac{s^2}{3})\big]
=-\sigma^2e^{-s}\,\tfrac{s(1+s)}{3}$, and $\frac{ds}{d\log\ell}=-s$; multiply.
Sample paths are twice differentiable. This kernel is usually a more realistic model than the squared exponential
for economic and financial series.

\paragraph{Periodic} (\code{PeriodicKernel}, $\sigma^2,\ell,p$): with $u=\pi\tau/p$,
\[
  k=\sigma^2\exp\!\Big(-\frac{2\sin^2u}{\ell^2}\Big),\quad
  \frac{\partial k}{\partial\log\sigma^2}=k,\quad
  \frac{\partial k}{\partial\log\ell}=k\,\frac{4\sin^2u}{\ell^2},\quad
  \frac{\partial k}{\partial\log p}=k\,\frac{2u\sin 2u}{\ell^2}.
\]
\emph{Derivation of the last one.} $\frac{du}{d\log p}=-u$ and
$\frac{d\sin^2u}{du}=\sin2u$, so $\frac{d\sin^2u}{d\log p}=-u\sin2u$. Then
multiply by $-2k/\ell^2$.

\paragraph{Sum and product} (\code{SumKernel}, \code{ProductKernel}): for
$k=k_1+k_2$ the gradient is the concatenation $(\nabla k_1,\nabla k_2)$. For
$k=k_1k_2$ the product rule gives $(k_2\nabla k_1,\ k_1\nabla k_2)$.
Hyperparameters are ordered as those of $k_1$ followed by those of $k_2$, recursively.

\subsection{Validity}

\begin{theorem}[Schur product theorem]\label{thm:schur}
If $A,B\in\R^{m\times m}$ are positive semi-definite, so is their elementwise
product $A\circ B$.
\end{theorem}
\begin{proof}
Write $B=\sum_k\lambda_ku_ku_k^\T$ with $\lambda_k\ge0$. Then $(A\circ B)_{ij}=\sum_k\lambda_k(u_k)_iA_{ij}(u_k)_j$,
so for any $x$,
$x^\T(A\circ B)x=\sum_k\lambda_k\,(x\circ u_k)^\T A\,(x\circ u_k)\ge0$.
\end{proof}

\begin{corollary}\label{cor:closure}
If $k_1,k_2$ are valid kernels, so are $k_1+k_2$, $k_1k_2$, $c\,k_1$ for $c>0$,
$k_1(\phi(t),\phi(t'))$ for any map $\phi$, and pointwise limits of valid kernels.
\end{corollary}
\begin{proof}
Gram matrices of $k_1+k_2$ are sums of positive semi-definite matrices. Those of $k_1k_2$ are Schur
products (Theorem~\ref{thm:schur}). A warped kernel's Gram matrix on $t_1,\dots,t_m$ is $k_1$'s
Gram matrix on $\phi(t_1),\dots,\phi(t_m)$. The positive semi-definite cone is closed, which covers limits.
\end{proof}

\begin{proposition}\label{prop:se-valid}
The squared exponential kernel is positive semi-definite, on $\R$ and on $\R^d$ with
$\norm{t-t'}$ in place of $|t-t'|$.
\end{proposition}
\begin{proof}
On $\R$, write
$e^{-(t-t')^2/2\ell^2} = g(t)\,g(t')\,e^{tt'/\ell^2}$ with $g(t)=e^{-t^2/2\ell^2}$. The
power series $e^{tt'/\ell^2}=\sum_{n\ge0}(tt')^n/(\ell^{2n}n!)$ is a limit of sums of
the rank-one kernels $t^nt'^n$, which are positive semi-definite. Multiplying by
$g(t)g(t')$ (a diagonal congruence $DKD$) preserves positive semi-definiteness, and so does scaling by $\sigma^2>0$. On
$\R^d$, $e^{-\norm{t-t'}^2/2\ell^2}=\prod_{i}e^{-(t_i-t_i')^2/2\ell^2}$ is a product of
valid kernels (Corollary~\ref{cor:closure}).
\end{proof}

\begin{proposition}\label{prop:periodic-valid}
The periodic kernel is the squared exponential kernel on $\R^2$, with the same $\ell$, applied
to the warping $\phi(t)=(\cos(2\pi t/p),\sin(2\pi t/p))$. It is therefore positive semi-definite.
\end{proposition}
\begin{proof}
$\norm{\phi(t)-\phi(t')}^2 = 2-2\cos(2\pi(t-t')/p) = 4\sin^2(\pi(t-t')/p)$, so
$\exp(-\norm{\phi(t)-\phi(t')}^2/2\ell^2)=\exp(-2\sin^2(\pi\tau/p)/\ell^2)$. Conclude with
Proposition~\ref{prop:se-valid} and Corollary~\ref{cor:closure}.
\end{proof}

The Matérn 5/2 kernel is positive semi-definite by Bochner's theorem. Its spectral
density is proportional to $(5/\ell^2+\omega^2)^{-3}$, which is non-negative and
integrable~\cite{rasmussen2006}. With $\sigma_n^2>0$, $K=K_f+\sigma_n^2I$ is positive
\emph{definite} for every kernel above. Its Cholesky factor therefore exists in exact
arithmetic. The floor \code{min\_noise\_variance} keeps it computable in floating
point.

%==============================================================================
\section{Posterior prediction}

\begin{theorem}[Predictive distribution]\label{thm:predictive}
Let $t_*$ be a vector of $m$ new times, $K_* = k(t,t_*)\in\R^{T\times m}$ and
$K_{**}=k(t_*,t_*)$. Under \eqref{eq:model}, with $K=K_f+\sigma_n^2I$,
\[
  f(t_*)\mid z\sim\N\big(K_*^\T K^{-1}z,\ K_{**}-K_*^\T K^{-1}K_*\big),\qquad
  z_*\mid z\sim\N\big(K_*^\T K^{-1}z,\ K_{**}-K_*^\T K^{-1}K_*+\sigma_n^2I\big).
\]
\end{theorem}
\begin{proof}
$(z,f(t_*))$ is jointly Gaussian with zero mean, $\Cov(z)=K$, $\Cov(f(t_*))=K_{**}$
and $\Cov(z,f(t_*))=\Cov(f(t),f(t_*))=K_*$, because the noise is independent of $f$. The first claim is the
conditioning lemma (proved in the Kalman filter note, Lemma~2.1). A new observation
$z_*=f(t_*)+\varepsilon_*$ adds independent noise, which adds $\sigma_n^2I$ to the covariance.
\end{proof}

The predictive mean $K_*^\T\alpha$, with $\alpha=K^{-1}z$, is a weighted sum of kernel
functions centred at the data. The variance is the prior variance minus what the
data explain. \code{GaussianProcess::predict} returns, for each $t_*$, the
marginal mean and variance of $y_*=m+s\,z_*$ (noise included), and
\code{forecast(H)} calls it with $t_*=T,\dots,T+H-1$.

\subsection{Long-horizon behavior}

\begin{proposition}[Mean reversion]\label{prop:reversion}
If $k(\tau)\to0$ as $|\tau|\to\infty$ (squared exponential, Matérn), then as
$t_*\to\infty$ the predictive mean of $y_*$ tends to $m$ and its variance to
$s^2(k(0)+\sigma_n^2)$.
\end{proposition}
\begin{proof}
Each entry of $K_*$ is $k(t_i-t_*)\to0$, while $K^{-1}$ and $z$ are fixed. So the mean
$K_*^\T K^{-1}z\to0$ and the variance $k(0)-K_*^\T K^{-1}K_*+\sigma_n^2\to k(0)+\sigma_n^2$.
Transform back with $y_*=m+sz_*$.
\end{proof}

\begin{proposition}[Periodic extrapolation]\label{prop:periodic}
For the periodic kernel with period $p$, the predictive mean and variance at
$t_*+p$ equal those at $t_*$.
\end{proposition}
\begin{proof}
$k(t_i,t_*+p)=k(t_i,t_*)$ and $k(t_*+p,t_*+p)=k(t_*,t_*)$, since $\sin^2$ has period $\pi$. Hence $K_*$ and
$K_{**}$ are unchanged, and so is everything in Theorem~\ref{thm:predictive}.
\end{proof}

So a decaying kernel is a local model. Its forecasts fade towards the mean
on the time scale of the length scale, and its intervals widen to the
unconditional spread. Long-range structure must be built into the kernel: a
periodic component (or a sum with one) repeats seasonality indefinitely, and a
product of periodic and decaying kernels lets the pattern drift.

%==============================================================================
\section{Marginal likelihood}\label{sec:ml}

Integrating out $f$, $z\sim\N(0,K(\theta))$, so the log marginal likelihood is
\begin{equation}\label{eq:lml}
  \log p(z\mid\theta) = -\tfrac12\, z^\T K^{-1}z \;-\; \tfrac12\log\det K \;-\; \tfrac T2\log2\pi .
\end{equation}
The first term rewards fitting the data. The second penalizes complexity: a
flexible kernel (short length scales, small noise) can explain many different
data sets, so it spreads its probability mass thinly and pays through a larger
$\log\det K$. The trade-off is automatic, an Occam's razor that needs no
validation set.

\begin{lemma}\label{lem:matrix-calculus}
For a smooth family of invertible symmetric matrices $K(\theta)$, with $K'=\partial K/\partial\theta$:
(i) $\partial K^{-1}/\partial\theta = -K^{-1}K'K^{-1}$;
(ii) $\partial\log\det K/\partial\theta = \tr(K^{-1}K')$ when $K\succ0$.
\end{lemma}
\begin{proof}
(i) Differentiate $KK^{-1}=I$: $K'K^{-1}+K\,\partial K^{-1}=0$.
(ii) $\det(K+\epsilon K')=\det K\,\det(I+\epsilon K^{-1}K')$, and
$\det(I+\epsilon A)=\prod_i(1+\epsilon\lambda_i(A))=1+\epsilon\tr A+O(\epsilon^2)$.
Hence $\log\det(K+\epsilon K')=\log\det K+\epsilon\tr(K^{-1}K')+O(\epsilon^2)$.
\end{proof}

\begin{theorem}[Gradient of the log marginal likelihood]\label{thm:gradient}
With $\alpha=K^{-1}z$,
\[
  \frac{\partial\log p(z\mid\theta)}{\partial\theta_j}
  = \tfrac12\,\alpha^\T\frac{\partial K}{\partial\theta_j}\alpha - \tfrac12\tr\!\Big(K^{-1}\frac{\partial K}{\partial\theta_j}\Big)
  = \tfrac12\tr\!\Big((\alpha\alpha^\T-K^{-1})\,\frac{\partial K}{\partial\theta_j}\Big).
\]
\end{theorem}
\begin{proof}
Differentiate \eqref{eq:lml} with Lemma~\ref{lem:matrix-calculus}:
$\partial(z^\T K^{-1}z)=-z^\T K^{-1}K'K^{-1}z=-\alpha^\T K'\alpha$ and
$\partial\log\det K=\tr(K^{-1}K')$. Finally, $\alpha^\T K'\alpha=\tr(\alpha\alpha^\T K')$.
\end{proof}

For a kernel hyperparameter, $\partial K/\partial\log\theta_j$ is the matrix of
\code{gradient()} entries. For the noise, $\partial K/\partial\log\sigma_n^2=\sigma_n^2I$, so
$\partial\log p/\partial\log\sigma_n^2=\tfrac12\sigma_n^2\tr(\alpha\alpha^\T-K^{-1})$.

\subsection{Normalization}\label{sec:normalization}

\begin{proposition}\label{prop:normalization}
If $z=(y-m\mathbf1)/s$ with fixed $m,s$, then $\log p(y)=\log p(z)-T\log s$.
\end{proposition}
\begin{proof}
The map $y\mapsto z$ is affine with Jacobian $s^{-T}$, so
$p_y(y)=p_z(z)\,s^{-T}$.
\end{proof}

\code{nabia} fits on $z$ when \code{normalize} is true (the default), so that
unit variances and noise levels around $10^{-2}$ are sensible starting points
whatever the units of $y$. The constant $-T\log s$ does not depend on $\theta$,
so the optimum is unchanged. \code{log\_marginal\_likelihood()} reports $\log p(y)$
in the units of $y$, using Proposition~\ref{prop:normalization}. Fitted kernel
variances and the noise variance are in normalized units: multiply by $s^2$ for
the units of $y$. The mean $m$ is the sample mean, plugged in rather than estimated
jointly. This is standard and harmless when the series has no strong trend.

%==============================================================================
\section{Computation}

\paragraph{Cholesky factorization.} Write $K=LL^\T$ with $L$ lower triangular (cost $T^3/3$). Then:
\begin{itemize}
  \item $\alpha=L^{-\T}(L^{-1}z)$, two triangular solves costing $O(T^2)$;
  \item $\log\det K=2\sum_i\log L_{ii}$, since $\det K=\det(L)^2=\prod_iL_{ii}^2$;
  \item for predictions, $v=L^{-1}K_*$ gives $K_*^\T K^{-1}K_*=v^\T v$, so the
        predictive variances are $\diag(K_{**})-\sum_i v_{i\cdot}^2$. This costs $O(T^2)$ per
        new time and never forms $K^{-1}$.
\end{itemize}
Subtracting $v^\T v$ can produce tiny negative values in floating point;
\code{predict} clamps the latent variance at $0$ before adding the noise.

\paragraph{Gradient.} Theorem~\ref{thm:gradient} needs $W=\alpha\alpha^\T-K^{-1}$.
\code{nabia} forms $K^{-1}$ from the Cholesky factor ($O(T^3)$) and accumulates
$-\tfrac12\sum_{ij}W_{ij}\,\partial k(t_i,t_j)/\partial\log\theta$ for all
hyperparameters in one pass over the $T^2$ pairs. One evaluation of the objective
and its gradient therefore costs $O(T^3+T^2n_\theta)$.

\begin{center}
\begin{tabular}{@{}lll@{}}
\toprule
Operation & Cost & Memory \\
\midrule
Gram matrix & $O(T^2)$ kernel evaluations & $O(T^2)$ \\
Cholesky, $\alpha$, $\log\det$ & $O(T^3)$ & $O(T^2)$ \\
Likelihood gradient & $O(T^3+T^2n_\theta)$ & $O(T^2)$ \\
Prediction at $m$ times & $O(T^2m)$ & $O(Tm)$ \\
\bottomrule
\end{tabular}
\end{center}
In practice this means a few thousand observations at most. Beyond that,
state-space representations of Matérn kernels (which reduce to Kalman filtering) or sparse
approximations are needed.

%==============================================================================
\section{Hyperparameter optimization}\label{sec:optimization}

The hyperparameters are
$\theta=(\log\theta_1,\dots,\log\theta_{n_k},\log\sigma_n^2)$, without the last entry when
\code{optimize\_noise} is false. \code{fit} minimizes $-\log p(z\mid\theta)$ with
\code{LBFGS} (tolerance $10^{-5}$ on the gradient, at most \code{max\_iter}
iterations), starting from the kernel and noise given in the parameters.
\begin{itemize}
  \item \emph{Why log-parameters.} Positivity becomes unconstrained, and the
        objective is much better conditioned. Length scales and variances matter on a
        relative scale, so a unit step in $\log\ell$ means the same thing for $\ell=0.1$ and $\ell=100$.
  \item \emph{Safeguards.} A trial with $|\theta_j|>20$ (a factor $e^{20}\approx5\cdot10^8$), or
        whose $K$ fails the Cholesky factorization, returns $+\infty$. The L-BFGS line search
        rejects such steps and backtracks.
  \item \emph{Noise floor.} The noise is clamped at \code{min\_noise\_variance}. In
        the clamped region the true gradient in $\log\sigma_n^2$ is zero, but the formula
        above is not. This only matters if the optimum lies at the floor, which
        signals noiseless data.
  \item \emph{Local optima.} The marginal likelihood is not concave and can be
        multimodal: for example, a short-length-scale explanation that fits the noise against a long-length-scale one with
        more noise, or periods that are multiples of each other. Start from plausible values,
        such as a period near the known seasonality. The \code{nabia} tests start a periodic kernel at
        $p=11.5$ and recover $p=12.0$. Restarting from several initial points is a cheap
        safeguard.
  \item The parameter struct keeps the initial kernel. The fitted copy is
        available through \code{kernel()}, and \code{noise\_variance()} returns the fitted noise.
\end{itemize}

%==============================================================================
\section{Implementation}

\subsection{Interface}
\begin{lstlisting}
auto p = std::make_unique<GaussianProcessParams>();
p->kernel = std::make_unique<SumKernel>(
    std::make_unique<Matern52Kernel>(1.0, 20.0),           // slow trend
    std::make_unique<PeriodicKernel>(0.5, 1.0, 12.0));     // seasonality
p->noise_variance = 1e-2;                                   // initial value
GaussianProcess gp(std::move(p));
gp.fit(y);                                // equally spaced, finite values
Forecast fc = gp.forecast(24);            // t = T .. T+23, noise included
Forecast at = gp.predict(times);          // any times, e.g. between observations
double lml = gp.log_marginal_likelihood();
\end{lstlisting}

\subsection{Algorithm}
\begin{algorithm}[ht]
\caption{\code{GaussianProcess::fit} and \code{predict}}\label{alg:gp}
\begin{algorithmic}[1]
\Require $y_0,\dots,y_{T-1}$ finite, $T\ge2$; kernel $k$; noise $\sigma_n^2$
\State $t\gets(0,\dots,T-1)$;\; $m\gets\bar y$;\; $s\gets$ sample sd (or 1);\; $z\gets(y-m)/s$
\If{\code{optimize\_hyperparameters}}
  \State $\theta^\star\gets$ L-BFGS minimizer of $-\log p(z\mid\theta)$, using Theorem~\ref{thm:gradient} for the gradient
  \State set $k$ and $\sigma_n^2$ from $\theta^\star$
\EndIf
\State $L\gets\mathrm{chol}(K_f+\sigma_n^2I)$;\; $\alpha\gets L^{-\T}L^{-1}z$ \Comment{stored for prediction}
\Statex \textbf{predict}$(t_*)$:
\State $K_*\gets k(t,t_*)$;\; $V\gets L^{-1}K_*$
\State $\mu\gets m+s\,K_*^\T\alpha$;\; $v\gets s^2\big(\max(\diag K_{**}-\mathrm{colsum}(V\circ V),0)+\sigma_n^2\big)$
\end{algorithmic}
\end{algorithm}

\paragraph{Validation.} The constructor requires a kernel and positive noise
variances. \code{fit} requires at least two finite observations, so the GP does not
accept missing values; use the Kalman filter for gappy series. Kernels
reject non-positive hyperparameters. A final Cholesky failure, which only happens if the
noise is too small for the kernel, throws \code{std::invalid\_argument}.

\paragraph{Correspondence with the code.}
\begin{center}
\begin{tabular}{@{}ll@{}}
\toprule
Mathematics & \code{GaussianProcess} / kernels \\
\midrule
$k(t,t')$, $\partial k/\partial\log\theta$ & \code{IKernel::operator()}, \code{gradient()} \\
$\log\theta$ & \code{log\_hyperparameters()}, \code{set\_log\_hyperparameters()} \\
$K_f$ & \code{IKernel::matrix(t, t)} \\
$m,\ s$ & private \code{offset\_}, \code{scale\_} \\
$L,\ \alpha$ & private \code{chol\_}, \code{alpha\_} \\
$\log p(y)$ & \code{log\_marginal\_likelihood()} \\
$m+sK_f\alpha$ & \code{fitted\_values()} \\
\bottomrule
\end{tabular}
\end{center}

\paragraph{Tests.} \code{tests/test\_forecasters.cpp} checks:
\begin{itemize}
  \item every kernel's gradient against central finite differences, plus symmetry;
  \item interpolation of a smooth function, with error below $10^{-3}$ between the data points;
  \item that hyperparameter fitting raises the log marginal likelihood by more than 10 nats
        and finds a plausible length scale;
  \item Proposition~\ref{prop:periodic}: a periodic kernel recovers period $12.0$,
        extrapolates 24 steps to within $0.1$, and the true values fall inside the 95\% intervals;
  \item Proposition~\ref{prop:reversion}: at $t_*=1000$ the forecast equals $m$
        and its variance equals $s^2(k(0)+\sigma_n^2)$, both to $10^{-12}$;
  \item rejection of invalid input.
\end{itemize}

\subsection{Limitations}
Cost is cubic in $T$. Inputs are equally spaced time indices, although \code{predict} accepts
any times. The mean is constant, the noise Gaussian and homoscedastic, and the
kernels stationary. There are no automatic restarts in the hyperparameter
search.

%==============================================================================
\begin{thebibliography}{9}
\bibitem{rasmussen2006} C.~E. Rasmussen and C.~K.~I. Williams, \emph{Gaussian Processes for Machine Learning},
  MIT Press, 2006.
\bibitem{roberts2013} S. Roberts, M. Osborne, M. Ebden, S. Reece, N. Gibson and S. Aigrain,
  Gaussian processes for time-series modelling, \emph{Phil. Trans. R. Soc. A} 371 (2013) 20110550.
\bibitem{mackay1998} D.~J.~C. MacKay, Introduction to Gaussian processes, in C.~M. Bishop (ed.),
  \emph{Neural Networks and Machine Learning}, Springer, 1998, 133--165.
\bibitem{horn2013} R.~A. Horn and C.~R. Johnson, \emph{Matrix Analysis}, 2nd ed., Cambridge University Press, 2013.
\bibitem{hartikainen2010} J. Hartikainen and S. S\"arkk\"a, Kalman filtering and smoothing solutions to temporal
  Gaussian process regression models, \emph{IEEE MLSP} (2010) 379--384.
\end{thebibliography}

\end{document}
